% Options for packages loaded elsewhere
\PassOptionsToPackage{unicode}{hyperref}
\PassOptionsToPackage{hyphens}{url}
%
\documentclass[
]{article}
\usepackage{amsmath,amssymb}
\usepackage{iftex}
\ifPDFTeX
  \usepackage[T1]{fontenc}
  \usepackage[utf8]{inputenc}
  \usepackage{textcomp} % provide euro and other symbols
\else % if luatex or xetex
  \usepackage{unicode-math} % this also loads fontspec
  \defaultfontfeatures{Scale=MatchLowercase}
  \defaultfontfeatures[\rmfamily]{Ligatures=TeX,Scale=1}
\fi
\usepackage{lmodern}
\ifPDFTeX\else
  % xetex/luatex font selection
\fi
% Use upquote if available, for straight quotes in verbatim environments
\IfFileExists{upquote.sty}{\usepackage{upquote}}{}
\IfFileExists{microtype.sty}{% use microtype if available
  \usepackage[]{microtype}
  \UseMicrotypeSet[protrusion]{basicmath} % disable protrusion for tt fonts
}{}
\makeatletter
\@ifundefined{KOMAClassName}{% if non-KOMA class
  \IfFileExists{parskip.sty}{%
    \usepackage{parskip}
  }{% else
    \setlength{\parindent}{0pt}
    \setlength{\parskip}{6pt plus 2pt minus 1pt}}
}{% if KOMA class
  \KOMAoptions{parskip=half}}
\makeatother
\usepackage{xcolor}
\usepackage{longtable,booktabs,array}
\usepackage{calc} % for calculating minipage widths
% Correct order of tables after \paragraph or \subparagraph
\usepackage{etoolbox}
\makeatletter
\patchcmd\longtable{\par}{\if@noskipsec\mbox{}\fi\par}{}{}
\makeatother
% Allow footnotes in longtable head/foot
\IfFileExists{footnotehyper.sty}{\usepackage{footnotehyper}}{\usepackage{footnote}}
\makesavenoteenv{longtable}
\usepackage{graphicx}
\makeatletter
\def\maxwidth{\ifdim\Gin@nat@width>\linewidth\linewidth\else\Gin@nat@width\fi}
\def\maxheight{\ifdim\Gin@nat@height>\textheight\textheight\else\Gin@nat@height\fi}
\makeatother
% Scale images if necessary, so that they will not overflow the page
% margins by default, and it is still possible to overwrite the defaults
% using explicit options in \includegraphics[width, height, ...]{}
\setkeys{Gin}{width=\maxwidth,height=\maxheight,keepaspectratio}
% Set default figure placement to htbp
\makeatletter
\def\fps@figure{htbp}
\makeatother
\setlength{\emergencystretch}{3em} % prevent overfull lines
\providecommand{\tightlist}{%
  \setlength{\itemsep}{0pt}\setlength{\parskip}{0pt}}
\setcounter{secnumdepth}{-\maxdimen} % remove section numbering
\ifLuaTeX
  \usepackage{selnolig}  % disable illegal ligatures
\fi
\IfFileExists{bookmark.sty}{\usepackage{bookmark}}{\usepackage{hyperref}}
\IfFileExists{xurl.sty}{\usepackage{xurl}}{} % add URL line breaks if available
\urlstyle{same}
\hypersetup{
  hidelinks,
  pdfcreator={LaTeX via pandoc}}

\author{}
\date{}

\begin{document}

\hypertarget{analyse-statistique-inter-canal-et-indicateurs-de-dispersion-locale}{%
\section{Analyse Statistique Inter-Canal et Indicateurs de Dispersion
Locale}\label{analyse-statistique-inter-canal-et-indicateurs-de-dispersion-locale}}

L\textquotesingle étude expérimentale de la répartition des couples de
candidats le long des axes de la structure modulaire \(\mathbb{R}_{30}\)
met en évidence des propriétés géométriques et de régularité
remarquables. En isolant le comportement local des flux au moyen
d\textquotesingle analyses par fenêtres glissantes et de calculs de
variance, il devient possible de quantifier la robustesse du modèle face
aux mécanismes d\textquotesingle élimination du crible.

\hypertarget{1-description-du-protocole-expuxe9rimental}{%
\subsection{1. Description du Protocole
Expérimental}\label{1-description-du-protocole-expuxe9rimental}}

Le modèle segmente l\textquotesingle espace des solutions paires en
canaux distincts fondés sur le groupe multiplicatif
\((\mathbb{Z}/30\mathbb{Z})^\times\). Pour un entier pair \(N\) donné,
les candidats \((p, q)\) respectant la relation \(p + q = N\) sont
regroupés selon leurs classes de congruence résiduelles mod 30.

\hypertarget{figure-1--matrice-dactivation-combinatoire-modulo-30}{%
\subsubsection{Figure 1 : Matrice d\textquotesingle Activation
Combinatoire Modulo
30}\label{figure-1--matrice-dactivation-combinatoire-modulo-30}}

\includegraphics{/home/mycmon/Documents/7. TH1-TH11-C1-C5/TH12/fig1_activation_matrix.png}

Cette matrice cartographie l\textquotesingle activation des classes
résiduelles finies et montre comment les interactions multiplicatives
s\textquotesingle organisent pour un \(N\) donné.

\begin{itemize}
\item
  XSG = Nombre Non Sophie Germain
\item
  SG = Nombre Sophie Germain
\end{itemize}

Afin de mesurer finement la dispersion et la régularité locale à
l\textquotesingle intérieur de ces canaux, l\textquotesingle espace de
chaque canal admissible est partitionné en \(B\) sous-blocs consécutifs
(par exemple, 8 blocs successifs contenant chacun un échantillon fixe de
1000 paires de candidats). Un processus de crible local de niveau
\(z = \sqrt{N}\) est appliqué à chaque bloc afin
d\textquotesingle évaluer le taux de survie des couples.

\hypertarget{figure-2--projection-circulaire-du-cercle-de-phase--ux1d445-30}{%
\subsubsection{\texorpdfstring{Figure 2 : Projection Circulaire du
Cercle de Phase
\(\mathbb{R}_{30}\)}{Figure 2 : Projection Circulaire du Cercle de Phase \textbackslash mathbb\{R\}\_\{30\}}}\label{figure-2--projection-circulaire-du-cercle-de-phase--ux1d445-30}}

\includegraphics{/home/mycmon/Documents/7. TH1-TH11-C1-C5/TH12/fig2_circle_R30.png}

Visualisation géométrique de l\textquotesingle espace des phases
\(\mathbb{R}_{30}\) mettant en évidence les symétries rigides des 8
classes de nombres premiers.

Les indicateurs statistiques suivants sont systématiquement calculés
pour chaque canal :

\begin{itemize}
\item
  \textbf{La moyenne locale (\(\mu\))} du taux de survie sur
  l\textquotesingle ensemble des blocs.
\item
  \textbf{L\textquotesingle écart-type (\(\sigma\))} et la
  \textbf{variance (\(\sigma^2\))} mesurant la fluctuation du taux
  d\textquotesingle un bloc à un autre.
\item
  \textbf{L\textquotesingle Indice de Dispersion (ou Facteur de Fano,
  \(I_D = \sigma^2 / \mu\))} caractérisant la nature de la distribution
  statistique sous-jacente.
\end{itemize}

\hypertarget{2-analyse-des-indicateurs-statistiques-et-interpruxe9tation}{%
\subsection{2. Analyse des Indicateurs Statistiques et
Interprétation}\label{2-analyse-des-indicateurs-statistiques-et-interpruxe9tation}}

Les données récoltées sur plusieurs ordres de grandeur (de \(10^8\) à
\(5 \times 10^9\)) révèlent trois caractéristiques majeures de
l\textquotesingle organisation modulaire.

\hypertarget{figure-3--enveloppe-spectrale-et-ondes-du-crible}{%
\subsubsection{Figure 3 : Enveloppe Spectrale et Ondes du
Crible}\label{figure-3--enveloppe-spectrale-et-ondes-du-crible}}

\includegraphics{/home/mycmon/Documents/7. TH1-TH11-C1-C5/TH12/fig3_spectral_envelope.png}

Illustration des forces d\textquotesingle exclusion du crible et de
leurs modulations harmoniques à grande échelle.

\hypertarget{a-la-sous-dispersion-extruxeame-et-lindice-de-fano}{%
\subsubsection{A. La Sous-Dispersion Extrême et l\textquotesingle Indice
de Fano}\label{a-la-sous-dispersion-extruxeame-et-lindice-de-fano}}

L\textquotesingle indicateur le plus significatif réside dans
l\textquotesingle affaiblissement systématique de
l\textquotesingle indice de dispersion (\(I_D\)). Dans un modèle
purement stochastique où la distribution des nombres premiers obéirait à
un processus de Poisson déconnecté de toute contrainte modulaire,
l\textquotesingle indice de dispersion fluctuerait autour de \(1\). Or,
les observations empiriques indiquent des valeurs uniformément basses,
chutant fréquemment sous le seuil de \(0.1\).

\hypertarget{graphique-b--analyse-de-lindice-de-dispersion-facteur-de-fano-par-canal}{%
\subsubsection{Graphique B : Analyse de l\textquotesingle Indice de
Dispersion (Facteur de Fano) par
Canal}\label{graphique-b--analyse-de-lindice-de-dispersion-facteur-de-fano-par-canal}}

Le graphique ci-dessous illustre ce fossé structurel : alors que le
hasard pur se situe sur la ligne rouge (\(1.0\)), les indices des canaux
réels restent écrasés au sol, confirmant une sous-dispersion
mathématique stricte.

L\textquotesingle alignement des ondes de crible sur la géométrie de la
roue de 30 engendre un effet de lissage : les éliminations de candidats
ne se regroupent pas de façon chaotique, interdisant localement
l\textquotesingle apparition de grands vides ou de « déserts » absolus
au sein du canal.

\includegraphics{/home/mycmon/Documents/7. TH1-TH11-C1-C5/TH12/monfette_dispersion_1000000008.png}

\hypertarget{b-luxe9quidistribution-des-canaux-simultanuxe9s}{%
\subsubsection{B. L\textquotesingle Équidistribution des Canaux
Simultanés}\label{b-luxe9quidistribution-des-canaux-simultanuxe9s}}

L\textquotesingle examen des moyennes de survie (\(\mu\)) entre les
différents canaux d\textquotesingle un même niveau \(N\) montre une
homogénéité remarquable. À l\textquotesingle échelle de \(N = 1\)
Milliard, les canaux configurés pour ce point se déploient avec des
moyennes de survie quasi-identiques : \textbf{7.12\%} pour le canal
\((1,17)\), \textbf{6.73\%} pour le canal \((7,11)\), et \textbf{7.02\%}
pour le canal \((19,29)\).

\hypertarget{graphique-a--uxe9volution-locale-du-taux-de-survie-par-bloc-consuxe9cutif}{%
\subsubsection{Graphique A : Évolution Locale du Taux de Survie par Bloc
Consécutif}\label{graphique-a--uxe9volution-locale-du-taux-de-survie-par-bloc-consuxe9cutif}}

On observe graphiquement la trajectoire lissée et parallèle de ces flux
à travers les blocs de candidats successifs, excluant toute anomalie
locale ou effondrement imprévu.

Cette observation valide la régularité du comportement des fonctions de
Dirichlet associées aux classes résiduelles de la roue de 30. Les forces
d\textquotesingle exclusion du crible s\textquotesingle exercent de
manière équitable, sans introduire de distorsion ou de biais majeur en
faveur d\textquotesingle un tunnel particulier.

\includegraphics{/home/mycmon/Documents/7. TH1-TH11-C1-C5/TH12/monfette_taux_survie_1000000008.png}

\hypertarget{figure-4--la-comuxe8te-de-goldbach-et-sa-distribution-macroscopique}{%
\subsubsection{Figure 4 : La Comète de Goldbach et sa Distribution
Macroscopique}\label{figure-4--la-comuxe8te-de-goldbach-et-sa-distribution-macroscopique}}

Perspective macroscopique classique de la distribution des couples,
servant de point de comparaison global avec la régularité microscopique
observée au sein des canaux.

\includegraphics{/home/mycmon/Documents/7. TH1-TH11-C1-C5/TH12/fig4_goldbach_comet.png}

\hypertarget{c-le-muxe9canisme-de-non-extinction-globale}{%
\subsubsection{C. Le Mécanisme de Non-Extinction
Globale}\label{c-le-muxe9canisme-de-non-extinction-globale}}

La trajectoire asymptotique analysée jusqu\textquotesingle à 5 Milliards
(\(N = 5\,000\,000\,030\)) confirme la robustesse géométrique du
dispositif. Face à l\textquotesingle augmentation drastique du nombre de
facteurs premiers testés (passant de 1 226 à 7 001 facteurs), la
déviation standard (\(\sigma\)) reste confinée aux alentours de
\textbf{1\%}.

\hypertarget{figure-5--diagramme-de-dualituxe9-de-phase-des-candidats}{%
\subsubsection{Figure 5 : Diagramme de Dualité de Phase des
Candidats}\label{figure-5--diagramme-de-dualituxe9-de-phase-des-candidats}}

Cartographie des symétries et des relations de dualité de phase entre
les composantes paires et impaires du modèle.

Même au sein des sous-blocs les plus exposés aux éliminations croisées,
le nombre absolu de paires survivantes reste largement supérieur à zéro.
La décroissance de la variance à mesure que \(N\) grandit rend la
probabilité d\textquotesingle une fluctuation simultanée sur
l\textquotesingle ensemble des canaux admissibles extrêmement
improbable.

\includegraphics{/home/mycmon/Documents/7. TH1-TH11-C1-C5/TH12/fig5_duality_diagram.png}

\hypertarget{figure-6--histogramme-des-indices-de-ruxe9gularituxe9-kappa--ux1d705-}{%
\subsubsection{\texorpdfstring{Figure 6 : Histogramme des Indices de
Régularité Kappa
(\(\kappa\))}{Figure 6 : Histogramme des Indices de Régularité Kappa (\textbackslash kappa)}}\label{figure-6--histogramme-des-indices-de-ruxe9gularituxe9-kappa--ux1d705-}}

Distribution statistique empirique validant les indices de régularité
mesurés sur l\textquotesingle ensemble de la campagne de calculs.

\includegraphics{/home/mycmon/Documents/7. TH1-TH11-C1-C5/TH12/fig6_kappa_histogram.png}

\hypertarget{3-synthuxe8se-des-observations-numuxe9riques}{%
\subsection{3. Synthèse des Observations
Numériques}\label{3-synthuxe8se-des-observations-numuxe9riques}}

Le tableau ci-dessous résume l\textquotesingle évolution des indicateurs
de dispersion et confirme la stabilisation des flux sur les
configurations testées :

\begin{longtable}[]{@{}lllllll@{}}
\toprule\noalign{}
\textbf{Valeur de N} & \textbf{Niveau du Crible (z)} & \textbf{Nombre de
Premiers} & \textbf{Canal Modulaire} & \textbf{Moyenne (μ)} &
\textbf{Écart-type (σ)} & \textbf{Indice de Dispersion (ID)} \\
\midrule\noalign{}
\endhead
\bottomrule\noalign{}
\endlastfoot
\textbf{\(100\,000\,018\)} & \(10\,000\) & 1 226 & \((11,17)\)
\((29,29)\) & \(6.812\%\) \(6.775\%\) & \(0.936\) \(0.772\) & \(0.1286\)
\textbf{\(0.0881\)} \\
\textbf{\(1\,000\,000\,008\)} & \(31\,622\) & 3 398 & \((1,17)\)
\((7,11)\) \((19,29)\) & \(7.125\%\) \(6.738\%\) \(7.025\%\) & \(1.036\)
\(0.583\) \(0.808\) & \(0.1508\) \textbf{\(0.0504\)} \(0.0931\) \\
\textbf{\(5\,000\,000\,030\)} & \(70\,710\) & 7 001 & \((1,19)\)
\((7,13)\) & \(5.225\%\) \(5.050\%\) & \(1.062\) \(0.886\) & \(0.2161\)
\(0.1554\) \\
\end{longtable}

\hypertarget{4-conclusion-provisoire}{%
\subsection{4. Conclusion Provisoire}\label{4-conclusion-provisoire}}

L\textquotesingle approche par indicateurs de dispersion montre que
l\textquotesingle extinction simultanée des canaux est contrecarrée par
deux facteurs empiriques : la très faible variance inter-blocs et la
valeur minimale de l\textquotesingle indice de Fano. Ces observations
suggèrent que la structure géométrique modulaire impose une contrainte
de régularité interne, garantissant le maintien de flux actifs de
candidats pour tout \(N\) pair échantillonné.

\hypertarget{figure-7--cube-topologique----3-uxd7-3-uxd7-3--des-configurations-minimales}{%
\subsubsection{\texorpdfstring{Figure 7 : Cube Topologique
\(3 \times 3 \times 3\) des Configurations
Minimales}{Figure 7 : Cube Topologique 3 \textbackslash times 3 \textbackslash times 3 des Configurations Minimales}}\label{figure-7--cube-topologique----3-uxd7-3-uxd7-3--des-configurations-minimales}}

Modélisation tridimensionnelle des contraintes de recouvrement non vides
(discutée en annexes ou en lien avec la formalisation Lean 4).

\includegraphics{/home/mycmon/Documents/7. TH1-TH11-C1-C5/TH12/fig7_cube_3x3x3.png}

\end{document}
